\documentclass[output=paper,colorlinks,citecolor=brown]{langscibook}
\ChapterDOI{10.5281/zenodo.21237001}
\title{To be or to become, that is the question: Word formation of the weak verbs formed from color adjectives in Old West Norse} 
\shorttitlerunninghead{Weak verbs formed from color adjectives in Old West Norse}

\author{Ramón Boldt\orcid{}\affiliation{Friedrich-Alexander-Universität
Erlangen-Nürnberg}}

\abstract{This article examines the semantic properties of a group of twelve Old West Norse deadjectival weak verbs, namely those derived from the eight Basic Color Terms that are attested in this language. The data is drawn from five major Old West Norse dictionaries. This way, it can be shown that the three class I \textit{jan}-verbs which are of Proto-Germanic origin share factitive semantics. The six genuinely Nordic class II \textit{nan}-verbs share anticausative-inchoative semantics. The three class II \textit{ōn}-verbs, however, are diverse in semantics and age. It can further be shown that the lack of context typical for dictionaries nurtures semantic imprecisions and inconsistencies, which can be met by investigating the original sources.}


\IfFileExists{../localcommands.tex}{
   \addbibresource{../localbibliography.bib}
   \usepackage{tabularx,multicol}
\usepackage{url}
\urlstyle{same}

 
\usepackage{langsci-optional}
\usepackage{langsci-lgr}
\usepackage{langsci-gb4e}
\usepackage[linguistics,edges]{forest}
\usepackage{tikz-qtree}
\usetikzlibrary{positioning}
\usetikzlibrary{patterns.meta}
\usepackage{multirow}
\usepackage{pgfplots}
\usepackage{setspace}
\usepackage{amsmath,stackengine}
% \usepackage{fontspec}
% \usepackage{tipa} % TIPA is banned
\usepackage{langsci-textipa}
\usepackage{langsci-branding}
\usepackage{longtable}
\usepackage{tabto}

\usepackage{enumitem}

%\usepackage{lscape} %Querformat chapter 4

   % \newcommand*{\orcid}{}


\makeatletter
\let\thetitle\@title
\let\theauthor\@author
\makeatother

\newcommand{\togglepaper}[1][0]{
   \bibliography{../localbibliography}
   \papernote{\scriptsize\normalfont
     \theauthor.
     \titleTemp.
     To appear in:
    Laura Grestenberger, Viktoria Reiter \& Melanie Malzahn (eds.).
    \emph{Deadjectival verb formation in Indo-European and beyond}.
     Berlin: Language Science Press. [preliminary page numbering]
   }
   \pagenumbering{roman}
   \setcounter{chapter}{#1}
   \addtocounter{chapter}{-1}
}

\newbool{bookcompile}
\booltrue{bookcompile}
\newcommand{\bookorchapter}[2]{\ifbool{bookcompile}{#1}{#2}}

%Sonderzeichen
%idg palatales k
\newcommand{\kinvbreve}{k\raisebox{0.6ex}{\hspace{-0.05em}\char"0311}}

\newcommand{\jac}{%
  \textit{\char"0237}% italic dotless j (ȷ)
  \hspace{-0.01em}% fine-tune horizontal spacing
  \raisebox{0.01ex}{\char"0301}% raised combining acute
}

%Baltic circumflex l
\newcommand\ltilded{%
\hspace{-0.2em}%
  \stackon[-1.5pt]{l}{\smash{\kern0.2em\~{}}}%
  \hspace{-0.2em}%
}

%Slav. wedge+ double acute
\newcommand{\HighH}[1]{%
    \stackon[-1.2ex]{#1}{\kern0.1em\H{}\kern-0.1em}%
}

%Lith dot tilde VICKY diese b.-sl. Sonderzeichen BRINGEN MICH UM!!
\newcommand{\tildedot}[1]{%
    \ensurestackMath{%
        \stackon[-0.27ex]{% Tilde layer 
            \stackon[-0.15ex]{% Dot layer 
                #1%
            }{\hspace{0.1em}\cdot\hspace{-0.13em}}%
        }{\hspace{0.1em}\text{\~{}}\hspace{-0.13em}%
    }%
}}

%Lith dot acute
\newcommand{\dotacute}[1]{%
    \ensurestackMath{%
        \stackon[-1.1ex]{% acute layer 
            \stackon[-0.15ex]{% Dot layer 
                #1%
            }{\hspace{0.1em}\cdot\hspace{-0.12em}}%
        }{\hspace{0.1em}\text{\'{}}\hspace{-0.12em}%
    }%
}}



% \setlength{\emergencystretch}{2pt} 

 \pgfplotsset{compat=1.18}


\defcitealias{AbB_1}{AbB 1}
\defcitealias{AbB_7}{AbB 7}
\defcitealias{ABL_1}{ABL 1}
\defcitealias{Vries1977}{AEW}
\defcitealias{ALEW}{ALEW}
\defcitealias{AnGr}{AnGr}
\defcitealias{ARM_1}{ARM 1}
\defcitealias{BT}{BT}
\defcitealias{black_concise_2000}{CDA}
\defcitealias{CHD}{CHD}
\defcitealias{CT28}{CT 28}
\defcitealias{CT34}{CT 34}
\defcitealias{CT39}{CT 39}
\defcitealias{Cleasby1874}{CV}
\defcitealias{Beekes2010}{EDG}
\defcitealias{eDiAna}{eDiAna}
\defcitealias{eDIL}{eDIL}
\defcitealias{EDL}{EDL}
\defcitealias{Kroonen2013}{EDPG}
\defcitealias{ESJS}{ESJS}
\defcitealias{EVP}{EVP}
\defcitealias{EWA1988}{EWA}
\defcitealias{EWAI}{EWAia I}
\defcitealias{EWAII}{EWAia II}
\defcitealias{EWGP}{EWGP}
\defcitealias{GDLI1961}{GDLI}
\defcitealias{GPC}{GPC}
\defcitealias{GRADIT1999}{GRADIT}
\defcitealias{HW}{HW}
\defcitealias{IEW}{IEW}
\defcitealias{KAH_2}{KAH 2}
\defcitealias{KAR_1}{KAR 1}
\defcitealias{KAR_2}{KAR 2}
\defcitealias{KEWA}{KEWA}
\defcitealias{Labat_TDP}{Labat TDP}
\defcitealias{lambert_BWL}{Lambert BWL}
\defcitealias{LEIA}{LEIA}
\defcitealias{Rix2001}{LIV\textsuperscript{2}}
\defcitealias{LIVaddenda}{LIV\textsuperscript{2} {A}ddenda}
\defcitealias{LKG1971}{LKG}
\defcitealias{LKZ}{LKŽ}
\defcitealias{LP}{LP}
\defcitealias{MhdGr}{MhdGr}
\defcitealias{MLLVG1959}{MLLVG}
\defcitealias{NIL}{NIL}
\defcitealias{OgnS}{OgnS}
\defcitealias{ONP}{ONP}
\defcitealias{RIMA_2}{RIMA 2}
\defcitealias{TIM_2}{TIM 2}
\defcitealias{WAP}{WAP}
\defcitealias{YOS_2}{YOS 2}






\newindex{wan}{wdx}{wnd}{Form index}
\newcommand{\iw}[2]{\index[wan]{#1!{#2}\xspace}}
\newcommand{\iwi}[3][]{\index[wan]{#2!{#3}\xspace#1}{#3}}
\newcommand{\iwit}[3][]{\rule{0pt}{0pt}#1\index[wan]{#2!{#3}}{#3}}


\newindex{van}{vdx}{vnd}{Form index}
\newindex{ban}{bdx}{bnd}{Book index}

% \newcommand{\iw}[3][]{\index[wan]{#2!\textit{#3}}}


% \providecommand{\iwi}[3][]{\iw{#3}\textit{#3}}


\newcommand{\indexnote}[1]{{\footnotesize\sffamily\raggedright\normalfont\color{gray}#1}}

   %% hyphenation points for line breaks
%% Normally, automatic hyphenation in LaTeX is very good
%% If a word is mis-hyphenated, add it to this file
%%
%% add information to TeX file before \begin{document} with:
%% %% hyphenation points for line breaks
%% Normally, automatic hyphenation in LaTeX is very good
%% If a word is mis-hyphenated, add it to this file
%%
%% add information to TeX file before \begin{document} with:
%% %% hyphenation points for line breaks
%% Normally, automatic hyphenation in LaTeX is very good
%% If a word is mis-hyphenated, add it to this file
%%
%% add information to TeX file before \begin{document} with:
%% \include{localhyphenation}
\hyphenation{
    par-a-digm non-de-adj-ect-iv-al -pre-sent -pre-sents for-ma-tion-al in-do-ger-ma-ni-sche in-do-ger-ma-ni-schen Rei-chert Rei-ter dia-chrony ety-mo-lo-gi-sches Ost-ger-ma-ni-schen ger-ma-ni-schen alt-in-di-schen mor-pho-lo-gi-scher Alt-indo-ari-sch-en lexico-graphy Schrij-ver
}
% aktuell in Bibliographie: stud-ies, histor-ical Soci-età
\hyphenation{
    par-a-digm non-de-adj-ect-iv-al -pre-sent -pre-sents for-ma-tion-al in-do-ger-ma-ni-sche in-do-ger-ma-ni-schen Rei-chert Rei-ter dia-chrony ety-mo-lo-gi-sches Ost-ger-ma-ni-schen ger-ma-ni-schen alt-in-di-schen mor-pho-lo-gi-scher Alt-indo-ari-sch-en lexico-graphy Schrij-ver
}
% aktuell in Bibliographie: stud-ies, histor-ical Soci-età
\hyphenation{
    par-a-digm non-de-adj-ect-iv-al -pre-sent -pre-sents for-ma-tion-al in-do-ger-ma-ni-sche in-do-ger-ma-ni-schen Rei-chert Rei-ter dia-chrony ety-mo-lo-gi-sches Ost-ger-ma-ni-schen ger-ma-ni-schen alt-in-di-schen mor-pho-lo-gi-scher Alt-indo-ari-sch-en lexico-graphy Schrij-ver
}
% aktuell in Bibliographie: stud-ies, histor-ical Soci-età
   \boolfalse{bookcompile}
   \togglepaper[2]%%chapternumber
}{}

\begin{document}
\SetupAffiliations{mark style=none}
\maketitle


\section{Introduction and aims}
\label{sec:OWNintro}

In what follows, I will present some initial considerations as part of an ongoing project on Old West Norse (OWN) verbal word formation. As the Vienna workshop focused on deadjectival verb formation,\is{deadjectival!verb} the topic of this paper is the formation of weak verbs to color adjectives. For this purpose, I have mainly evaluated some of the larger Old West Norse dictionaries (Section \ref{sec:OWNSources}). Underlining the shortcomings of this method in comparison to working with corpora with regard to issues of \textit{Aktionsart} (hence ``to be or to become" in the title) has also been a goal of the talk and, subsequently, of this article.

\section{Basic Color Terms}
\label{sec:OWNBCT}

With respect to color adjectives, it is not possible to cover all words for colors that are attested in Old West Norse texts for reasons of space. Therefore, the selection will be limited to the so-called Basic Color Terms,\is{Basic Color Term (BCT)} a term well established since \citealt{Berlin1969} (\textit{Basic Color Terms: Their Universality and Evolution}). In the most recent investigation, Jackson Crawford’s dissertation \textit{The Historical Development of Basic Color Terms in Old Norse-Icelandic} (\citealt{Crawford2014}), the author provides a selection of color terms that can be considered basic in Old West Norse.\footnote{For a similar approach, cf. \citet{Wolf2013}, \citet{Bruckmann2012}, \citet{Wolf2006a}. For the Non-Basic Color Terms in Old West Norse, cf. \citet{Wolf2015}.} Crawford himself refers to the latest comprehensive study of the topic, \citet{Biggam2012}. Here is not the place to discuss their concept of \isi{Basic Color Term (BCT)} in length. Instead, I will present the most important criteria according to \citet[125--153]{Crawford2014} and \citet[21--42]{Biggam2012} in order to motivate the focus on a certain class of color adjectives in this article. In principle, I follow both authors, but not always rigidly, and I deviate from their criteria when it serves my purpose of elucidating matters of verbal word formation.

\begin{enumerate}

    \item[a)] The color term in question must be monolexemic, i.e., a simplex. It must therefore not be deducible from the sum of its parts. In other words, it must be non-predictable. Most of the complex color terms in Old West Norse seem to be \textit{ad hoc} formations, attested only once or twice, e.g., \iwi{ON}{\textit{móbrúnn}} ‘dark brown’ or \iwi{ON}{\textit{iðjagrǿnn}} ‘green again (?)’. However, it can be argued that those compounds\is{compounding} could only have been formed on the assumption that readers must have been able to predict their meaning.\footnote{\citet[126]{Crawford2014} discusses the compound\is{compounding} \iwi{ON}{\textit{skolbrúnn}}, which he translates as “dirty water brown” or perhaps “dirty-water (eye)browed”. In light of similar compounds\is{compounding} like \iwi{ON}{\textit{léttbrúnn}} ‘light-browed (i.e., happy looking)’ and \iwi{ON}{\textit{hvítbrúnn}} ‘white-browed’, it seems to me that this can not be taken as a good example here.}

    \item[b)] The color term in question must not be a hyponym. \citet[59]{Crawford2014} is “hesitant” in ascribing basic status to OWN \iwi{ON}{\textit{brúnn}} ‘brown’. He rather considers it a contextually restricted hyponym for \iwi{ON}{\textit{svartr}} ‘black’, as it is used only to describe a certain dark hue of hair, especially horse hair. In other words, a horse with \textit{brúna} hair can also be described as having \textit{svarta} \iw{ON}{\textit{svartr}} hair. Nonetheless, I will make some brief remarks on OWN \iwi{ON}{\textit{brúnn}} here as well, because it is important for issues related to word formation.

    \item[c)] Related to b) is the idea that the use of a color term in question must not be restricted to a class of objects. A good example can be seen in German and English \textit{blond} \iw{German}{\textit{blond}} which can only be used to describe the human hair color.\footnote{In contrast to OWN \iwi{ON}{\textit{brúnn}} : \iwi{ON}{\textit{svartr}}, English and German \textit{blond} is usually not considered a hyponym of, say, \iwi{German}{\textit{gelb}} or \textit{yellow}, respectively.}

    \item[d)] BCTs \is{Basic Color Term (BCT)} are morphologically more productive\is{productivity} than other color terms. For example, in contemporary English, we see that only \textit{red}, \textit{black}, and \textit{white} can serve as bases for \isi{anticausative}-\isi{inchoative} derivations: \textit{redden}, \textit{blacken}, and \textit{whiten}. This could, as we will see, be true for Old West Norse as well in that at least one derivative is attested for every BCT.\is{Basic Color Term (BCT)}

    \item[e)] Finally, a lexeme in question must contain information about hue, that is “the spectrum of visible light, parts of which, according to their wavelength or frequency, are perceived by humans to differ from others” \citep[3]{Biggam2012}. Lexemes that only contain information on saturation or tone, that is “the admixture of white or black with a hue” (ibid.), are excluded. An example from OWN would be \iwi{ON}{\textit{døkkr}} ‘dark’, which is neutral in respect to hue. However, OWN \iwi{ON}{\textit{bleikr}}, often translated as ‘pale’ in the dictionaries or conceptualized as “macrocolor covering, at least partly, the category of pale of light colors” \citep[143]{Wolf2013}, is part of Crawford’s investigation as he favors it over \iwi{ON}{\textit{gulr}} as the BCT \is{Basic Color Term (BCT)} for ‘yellow’. I follow him in that regard.

\end{enumerate}

In light of what we have established so far, the following color terms will be considered basic in Old West Norse: \iwi{ON}{\textit{blár}} ‘blue, black’; \iwi{ON}{\textit{bleikr}} ‘yellow’; \iwi{ON}{\textit{brúnn}} ‘brown’; \iwi{ON}{\textit{grár}} ‘grey’; \iwi{ON}{\textit{grǿnn}} ‘green’; \iwi{ON}{\textit{hvítr}} ‘white’; \iwi{ON}{\textit{rauðr}} resp. \iwi{ON}{\textit{rjóðr}} ‘red’; \iwi{ON}{\textit{svartr}} ‘black’. Accordingly, the verbal word formation of those color adjectives will be the focus of this article.

\section{Sources}
\label{sec:OWNSources}

As mentioned in Section \ref{sec:OWNintro}, I focus on dictionaries here in order to highlight the shortcomings of this method in comparison to working with corpora, namely the sometimes imprecise determination of the semantics of the respective verbs. The following dictionaries were used:

\begin{itemize}
\item Cleasby’s and Vigfússon’s \textit{Old Norse to English Dictionary} (from now on \citetalias{Cleasby1874}),
\item Baetke’s \textit{Wörterbuch zur altnordischen Prosaliteratur} (\citetalias{WAP}), 
\item Fritzner’s \textit{Ordbog over det gamle norske Sprog} (\citetalias{OgnS}), 
\item and the not yet completed \textit{Dictionary of Old Norse Prose} (\citetalias{ONP}). 
\end{itemize}

The last one is certainly more than a mere dictionary in that it contextualizes the words in question and provides scans of the relevant editions of those texts in which the words are attested. Weak verbs from the poetic language recorded in the \textit{Lexicon Poeticum} (\citetalias{LP}) do not provide any substantial additional information regarding the questions discussed here. The \citetalias{LP} will therefore be quoted only occasionally.

\section{History of research}
\label{sec:OWNhistory}

In comparison to the other Old Germanic languages, we know relatively little about verbal word formation in Old (West) Norse, although researchers have long taken a keen interest in the language.\footnote{\label{FNOWNlit}For Gothic see \citet{MartiHeinzle2023}, \citet{GarciaGarcia2005}; for Old English see \citet{Bammesberger1965}; for Old High German see \citet{MartiHeinzle2019}, \citet{Riecke1996}; for Middle High German see \citet{Leipold2006}; for Early New High German see \citet{Habermann1994}; for Old Frisian see \citet{MartiHeinzle2014}; for a comparative study on Gothic and Continental West Germanic see \citet{Schwerdt2008}.} In the early Neogrammarian period, research focused mainly on formal aspects of verbal inflection. Noreen (\citeyear{AnGr}) – considered the standard reference grammar for Old West Norse until today – does not mention word formation at all. The same is true for other grammars and surveys from that era (\citealt{Wimmer1871}; \citealt{Brenner1882}; \citealt{Heusler1932}; \citealt{Krause1948}; \citealt{Andersen1954}). In contrast, \citet{Holthausen1895} seems to be the first to discuss the verbal word formation of Old West Norse. He does so on one page in a chapter entitled \textit{Bildungslehre} (131). Here, Holthausen is mostly concerned with the \textit{jan}-suffix that, according to him, primarily forms \isi{causative}s from strong ablauting verbs, for example, \iwi{ON}{\textit{føra}} ‘to lead, transport, move’ from \iwi{ON}{\textit{fara}} ‘to go, travel, fare’.\footnote{Still visible today in Nynorsk \textit{å føra} : \textit{å fara}\iw{ModNorw}{\textit{å føra} (Nynorsk)} \iw{ModNorw}{\textit{å fara} (Nynorsk)} or German \iwi{German}{\textit{führen}} : \iwi{German}{\textit{fahren}}.} Additionally, the class II \textit{ōn}-suffix forms denominatives,\is{denominative!verb} e.g., \iwi{ON}{\textit{tala}} ‘to speak’ from \iwi{ON}{\textit{tal}} ‘talk, conversation; speech, language’. Also, class II contains \textit{nan}-verbs of the former class IV, e.g., \iwi{ON}{\textit{vakna}} ‘wake up’ in contrast to \iwi{ON}{\textit{vaka}} ‘being awake, waking’. Holthausen describes those verbs as passive and \isi{inchoative}, a view still upheld by some researchers.\footnote{\label{FNOWNinchoativity}In light of more recent research, it is argued that the class-forming feature of the \textit{nan}-verbs is not inchoativity \is{inchoative} (i.e., the verb denotes the (gradual) beginning of a situation) and/or passivity, but anticausativity \is{anticausative} (i.e., the verb denotes a situation where the subject is non-\isi{agentive}, often relating to internal and external nature), cf. \citet{Scheungraber2014}; \citet{Ottosson2013}; also \citet{VillanuevaSvensson2011}. \citet[16--17]{Haspelmath1987} seems to be the first to explicitly refer to the Gothic \textit{nan}-verbs as \isi{anticausative}. \citegen{Gorbachov2007} study on the nasal presents\is{nasal present} in Germanic, Slavic, and Baltic does not make use of the term \textit{anticausative} \is{anticausative} but calls this class \textit{inchoative}.\is{inchoative} See the introduction of this volume on the different uses of these terms in the context of deadjectival verb formation\is{deadjectival!verb} in general.}

Class IV exists as such only in Gothic. In West Germanic, these verbs joined class II or class III (in Old High German). A few \textit{nan}-verbs may be attested in Proto-Norse, but their readings and interpretations must be considered uncertain.\footnote{Cf.\ PN \iwi{PN}{\textit{hleunan}*} ‘to become warm’ (cf. \citealt{Boldt2024}), \iwi{PN}{\textit{sufnan}*}  ‘to start moving’, \iwi{PN}{\textit{dwenan}*}  ‘to dwindle’.} Moreover, it is difficult to determine whether they formed a separate class.

The path of research that began with \citet{Holthausen1895} seems to have reached its climax with \citegen{Wessen1965} \textit{Svensk språkhistoria: Ordbildningslära}. Wessén was comprehensive in his analysis, writing for almost 20 pages on Norse verbal formation (90–107). However, we are still rather limited in what we know compared to other older Germanic languages (see fn. \ref{FNOWNlit}). Except for the field’s current focus on the former class IV verbs, the state of research seems to be condensed in \citegen[118--120]{Nedoma2010} \textit{Kleine Grammatik des Altisländischen}. According to Nedoma, the three weak classes of Old West Norse entail the following semantic properties: 

\begin{itemize}
    \item class I (\textit{i)jan}-verbs: deverbal \isi{causative}s,\is{deverbal/deverbative!verb} denominal \is{denominal!verb} \isi{factitive}s, and some \isi{intensive}-iteratives\is{iterative}
    \item class II \textit{ōn}-verbs: denominals \is{denominal!verb} of various kinds, as well as some \isi{intensive}-iteratives\is{iterative}
    \item class III \textit{ēn}-verbs: mainly \isi{intransitive} \isi{durative}s.
\end{itemize}

\section{Weak verbs attested in Old West Norse}
\label{sec:OWNdata}
\subsection{General remarks}
\label{sec:OWNgeneral}

In what follows, I will evaluate the semantics of the Old West Norse weak verbs formed to BCTs \is{Basic Color Term (BCT)} as given in the dictionaries listed above and examine whether their translations correspond to what is discussed in the reference grammars and handbooks. In OHG, we can witness a decreasing correlation between semantic classes and certain morphemes, a process coming to an end in MHG, where the suffixes merely function as purely morphological markers of inflection.\footnote{\citetalias[ III: 494--496]{MhdGr}; \citet[109--111]{Boldt2023}.}

I will present the BCTs \is{Basic Color Term (BCT)} (as mentioned in Section \ref{sec:OWNBCT}) one after another in alphabetical order. I will also present their meaning, cognates in the other Old Germanic languages, basic etymological information, and finally the weak verbs derived from them. I will then give the meaning of those verbs as posited in the dictionaries as well as representative examples of their actual occurrences. In most cases, the meanings provided by the dictionaries agree with one another.  However, there are some instances where the provided meanings are inconsistent. In examples of the latter, I will try to provide a solution based on the actual texts.

\subsection{\textit{blár} ‘blue, blue-black; black; dark’}
\label{sec:OWNblue}

The first BCT \is{Basic Color Term (BCT)} to be presented here is \iwi{ON}{\textit{blár}} ‘blue; blue-black; black; dark’.\footnote{For a detailed study on the use of \iwi{ON}{\textit{blár}}, also in comparison to \iwi{ON}{\textit{svartr}}, in Old Icelandic, cf. \citet{Wolf2006b}.} It is well attested in the other Old Germanic languages, cf. OE in \iwi{OE}{\textit{blǽ-hǽwen}} ‘of a blue hue, bluish’ and \iwi{OE}{\textit{blǽwen}} ‘light blue’, OHG \iwi{OHG}{\textit{blāo}}, OFris. \iwi{OFris}{\textit{blāu}}, OS \iwi{OS}{\textit{blāo}} < PG \iwi{PGmc}{*\textit{blēwa}-} ‘blue’, and connected to Lat. \textit{flāvus}\iw{Latin}{\textit{flāuus}} and OIr. \iwi{OIr}{\textit{blā}}, although the exact circumstances are not entirely clear (cf. \citetalias[ II 160]{EWA1988}; \citetalias[ 68]{Kroonen2013}; \citetalias[ 135]{EWGP}; \citealt[63]{Luhr2000}).

Regarding verbal derivations\is{derivation!verbal} of \iwi{ON}{\textit{blár}}, only one weak verb is attested in Germanic in general, and in OWN in particular. It is the verb \textit{blána}\iw{ON}{\textit{blána} (OWN)} that belongs to the former class IV, but synchronically to class II. It is also attested in OSwed. \iwi{OSwed}{\textit{blana}}. According to \citetalias{ONP}, the verb is attested at least 13 times in OWN prose. With regard to the meaning of this verb, the dictionaries agree: They all imply \isi{anticausative}-\isi{inchoative} semantics.\footnote{As stated in fn.\ \ref{FNOWNinchoativity}, in more recent research the \isi{anticausative} character of (former) \textit{nan}-verbs is favored over the earlier assumption of primarily \isi{inchoative} semantics. As this issue is not decisive here, I will stick with the term \isi{anticausative}-\isi{inchoative}, cf. \citet{VillanuevaSvensson2011}.} 

\begin{itemize}
\item \citetalias{ONP}: ‘to turn black-and-blue, take on a livid color’
\item \citetalias{Cleasby1874}: ‘to become black, livid’
\item \citetalias{WAP}: `dunkel werden, eine dunkle Farbe annehmen, blauschwarz werden'
\item \citetalias{OgnS}: `blive blaa' (`to become blue')
\end{itemize}

An example for this use of \textit{blána}\iw{ON}{\textit{blána} (OWN)} is given in (\ref{ex:blana}):

\begin{exe}
    \ex Ágrip af Noregs konunga sǫgum 5,6\footnote{All sources are cited according to the editions that are used and given in the \citetalias{ONP}.} \label{ex:blana}\\
    \gll \textit{Var} \textit{þá} \textit{hvata-t} \textit{bál-i} \textit{ok} \textit{hon} \textit{bren-d} {\textit{\textbf{blána-þi}} \iw{ON}{\textit{blána} (OWN)}} \textit{þó} \textit{áþr} \textit{all-r} \textit{líkam-inn} \textit{ok} \textit{ull-o} \textit{ór} \textit{orm-ar} \textit{ok} \textit{eþl-or ...}\\
    was then {prepare.hastily-\textsc{ptcp.prf.nom.sg}} pyre-\textsc{dat.sg} and she burn-\textsc{ptcp.prf.nom.sg} \textbf{turn.dark}-\textsc{3sg.pret} but before complete-\textsc{nom.sg} body-\textsc{nom.sg.det} and swarm-3sg.\textsc{pret} out worm-\textsc{nom.pl} and serpent-\textsc{nom.pl} \\
    \glt ‘Then a pyre was hastily prepared, and she was burnt, but before that, the body \textbf{turned dark} completely and there swarmed out worms and serpents, ...’\\
\end{exe} 

\subsection{\textit{bleikr} ‘pale; yellow’}
\label{sec:OWNyellow}

The second BCT \is{Basic Color Term (BCT)} in our list is \iwi{ON}{\textit{bleikr}} ‘pale; (light) yellow’. In light of OE \iwi{OE}{\textit{blác}}, \iwi{OE}{\textit{blǽce}}, OHG \iwi{OHG}{\textit{bleih}}, and OS \iwi{OS}{\textit{blēk}}, a preform PG \iwi{PGmc}{*\textit{blaika}-} must be reconstructed (\citetalias[ II 176]{EWA1988}; \citetalias[ 66]{Kroonen2013}; \citetalias[ 127]{EWGP}). 

For \iwi{ON}{\textit{bleikr}}, three verbal derivations\is{derivation!verbal} are attested, the first one of which is the class I verb \iwi{ON}{\textit{bleikja}}. It is also attested in OE \iwi{OE}{\textit{blǽcan}}, OHG \iwi{OHG}{\textit{bleichen}}, MLG \iwi{MLG}{\textit{blēken}}, and must therefore be reconstructed as PG \iwi{PGmc}{*\textit{blaikijan}-}. The dictionaries once again agree on the meaning:

\begin{itemize}
\item \citetalias{ONP}: ‘to bleach, make lighter’
\item \citetalias{Cleasby1874}: ‘to bleach’
\item \citetalias{WAP}: ‘bleichen’
\item \citetalias{OgnS}: ‘blege’ (‘to bleach’)
\end{itemize}

Compare example (\ref{ex:bleikja}): 
 
\begin{exe} 
    \ex Snorra Edda 130,38 \label{ex:bleikja}\\
    \gll \textit{Brynhild-r} \textit{ok} \textit{Gvðrvn} \textit{gengv} \textit{til} \textit{vat-z} \textit{at} {\textit{\textbf{bleik-ia}} \iw{ON}{\textit{bleikja}}} \textit{hadd-a} \textit{sina}\\
       Brunhild-\textsc{nom.sg} and Guðrun.\textsc{nom.sg} went to water-\textsc{gen.sg} {to} \textbf{bleach}-\textsc{inf} hair-\textsc{acc.pl} their\\
    \glt ‘Brunhild and Guðrun went to the water in order to \textbf{bleach} their hair.’
\end{exe} 



The second weak verb derived from \iwi{ON}{\textit{bleikr}} is OWN \textit{blikna}\iw{ON}{\textit{blikna} (OWN)}. Just like \textit{blána}\iw{ON}{\textit{blána} (OWN)}, it is a class II verb with \isi{anticausative}-\isi{inchoative} semantics, earlier associated with the fourth class. As with \textit{blána}\iw{ON}{\textit{blána} (OWN)}, it is attested only in OWN (at least 30 times according to \citetalias{ONP}, twice according to \citetalias{LP}).\footnote{Cf. Middle English \iwi{ME}{\textit{bliknen}} ‘to become pale’ is a loan from OWN showing \isi{anticausative}-\isi{inchoative} semantics as well.} In addition, \textit{blikna}\iw{ON}{\textit{blikna} (OWN)} is one of the zero-grade derivations from a strong verb well attested in OWN,\footnote{On that type, cf. \citet[344]{Ottosson2013} with literature.} namely from OWN \iwi{ON}{\textit{blíkja}} ‘to glisten, glint’ (< \iwi{PGmc}{*\textit{bleikan}-} with secondary palatalization, see \citetalias[ §362]{AnGr}). This makes it likely that this is a purely Norse innovation.

In this case, the dictionaries again agree on the semantics of the verb in question:

\begin{itemize}
\item \citetalias{ONP}: ‘to turn pale, turn sallow’
\item \citetalias{Cleasby1874}: ‘to become pale’
\item \citetalias{WAP}: ‘blass werden, erbleichen’
\item \citetalias{OgnS}: ‘blegne, blive hvid eller bleg’ (‘to become pale, become white or sallow’)
\end{itemize}
 
\begin{exe} 
    \ex Óláfs saga helga 173,13 \label{ex:blikna}\\
    \gll \textit{þa} {\textit{\textbf{blicna-ði}} \iw{ON}{\textit{blikna} (OWN)}} \textit{hann} \textit{oc} \textit{varð} \textit{faul-r} \textit{sem} \textit{nár}\\
       then \textbf{turn.sallow}-\textsc{3sg.pret} he and become.\textsc{3sg.pret} pale-\textsc{nom.sg} like corpse\\
    \glt ‘He then \textbf{turned sallow} and became pale like a corpse.’
\end{exe}

Apart from \textit{blikna}\iw{ON}{\textit{blikna} (OWN)}, we have another class II verb belonging to \iwi{ON}{\textit{bleikr}} in OWN, namely \iwi{ON}{\textit{blika}}. In contrast to \textit{blikna}\iw{ON}{\textit{blikna} (OWN)}, it is an old \textit{ōn}-verb, going back to PG \iwi{PGmc}{*\textit{blikōn}-}, as attested in OSwed. \iwi{OSwed}{\textit{blika}}, OE \iwi{OE}{\textit{blician}} (\citetalias[ 44]{Vries1977}), another zero-grade formation of \iwi{PGmc}{*\textit{bleikan}-}. In contrast to the aforementioned verbal formations, the dictionaries disagree with each other in this case:

\begin{itemize}
\item \citetalias{ONP}: ‘to glisten, glint, shine’
\item \citetalias{Cleasby1874}: ‘to gleam, twinkle’
\item \citetalias{WAP}: ‘blinken, funkeln’ but also ‘erglänzen, aufblitzen'
\item \citetalias{OgnS}: = \textit{blíkja} ‘shine, gleam’
\item \citetalias{LP}: ‘gleam, flash’
\end{itemize}

Already \citet[240]{KraheMeid1969} stated that the old deverbal \textit{ōn}-verbs\is{deverbal/deverbative!verb} are predominantly of \isi{intensive}-\isi{iterative} character, which would be confirmed by the derivation’s meaning ‘to glisten, glint’ with regard to a verbal base \iwi{PGmc}{*\textit{bleikan}-} ‘to gleam, shimmer’. The entry in \citetalias{WAP}, however, would suggest an additional \isi{inchoative} or rather \isi{semelfactive} meaning ‘to flash’. To address this problem, I have examined the actual attestations. According to \citetalias{ONP}, \iwi{ON}{\textit{blika}} is attested at least 10 times in the corpus of Old West Norse prose. Strikingly, in 9 out of those 10 attestations \iwi{ON}{\textit{blika}} refers to \iwi{ON}{\textit{skildir}} ‘shields’, and the combination of both might be regarded as a fixed expression.\footnote{A fixed order of the two components could also account for the verb fronting in example sentences (\ref{ex:blika3}) and (\ref{ex:blika4}). On the other hand, the three attestations of \iwi{ON}{\textit{blika}} given in \citetalias{LP} from the poetic texts are not combined with \iwi{ON}{\textit{skildir}}, but with \iwi{ON}{\textit{merki}} ‘banners’, \iwi{ON}{\textit{vôpn}} ‘weapons’, and \iwi{ON}{\textit{reiði}} ‘harness’.} In light of this information, a classification of \iwi{ON}{\textit{bleikr}} ‘pale; (light) yellow’ as BCT \is{Basic Color Term (BCT)} might seem questionable, because it contradicts the fact that a BCT must not be restricted to a class of objects. As I cannot discuss this matter in detail here, I will keep the focus on issues of verbal word formation and therefore quote two passages:
 
\begin{exe}
    \ex Njáls saga 210,27 \label{ex:blika1}\\
    \gll ... \textit{Skild-ir} {\textit{\textbf{blik-a}}...,} \textit{er} \textit{sól-in} \textit{skin-n} {\textit{á} ...}\\
      {} shield-\textsc{nom.pl} \textbf{blika}-\textsc{3pl.prs} which sun.\textsc{nom.sg-det} shine-\textsc{3sg.prs} on \\
    \glt ‘The shields \textbf{\textit{blika}},\iw{ON}{\textit{blika}} on which the sun shines ...’\\
    
    \ex Heiðarvíga saga 99,16 \label{ex:blika2}\\
    \gll ... \textit{skip} \textit{skrið-r}, \textit{skilld-ir} {\textit{\textbf{blik-a}}} \textit{sol-skin} \textit{sne} \textit{leg-ir} \\
     {}  ship.\textsc{nom.sg} glide-\textsc{3sg.prs} shield-\textsc{nom.pl} \textbf{blika}-\textsc{3pl.prs} sun-shine.\textsc{nom.sg} snow.\textsc{acc.sg} melt-\textsc{3sg.prs}\\
    \glt ‘[Whoever breaks an oath is outlawed as long as Christians go to church, fires blaze, fields green], ship glides, shields \textbf{\textit{blika}},\iw{ON}{\textit{blika}} sunshine melts snow.’
\end{exe}

Neither of the two text passages can help us in determining the semantic value of \iwi{ON}{\textit{blika}} as \isi{intensive} or \isi{semelfactive}, since both are decidedly ambiguous. 

The second text passage in (\ref{ex:blika2}) is cited here additionally in order to draw attention to the shortcomings of working with very limited contexts. The entry given in the \citetalias{ONP} consists of nothing more than \textit{skilldir blika}. Though this is more than what we get from the other dictionaries, it still can hardly be called context at all. Luckily, this problem is resolved by the \citetalias{ONP}, which provides detailed information on where to find the passage in question in the edition as well as a scan of that very passage.

The other passages are not helpful in further determining the semantic value of \iwi{ON}{\textit{blika}}. To illustrate this claim I will quote another two of them:
 
\begin{exe}
    \ex Egils saga 159,17 \label{ex:blika3}\\
    \gll \textit{Menn} \textit{sa} \textit{af} \textit{þinginu} \textit{at} \textit{flock-r} \textit{manna} \textit{reið} \textit{neðan} {...} \textit{ok} {\textit{\textbf{bliku-ðu}} \iw{ON}{\textit{blika}}} \textit{þar} \textit{skilld-er} \textit{við}\\
    one saw from thing.\textsc{dat.sg.det} at crowd-\textsc{nom.sg} man.\textsc{gen.sg} ride.\textsc{3sg.pret} down {} and \textbf{blika}-\textsc{3pl.pret} there shield-\textsc{nom.pl} with\\
    \glt ‘One saw from the Thing a crowd of men riding down ... and there their shields \textit{\textbf{blikuðu}} ...’\\

    \ex Bolla þáttr 299,18 \label{ex:blika4}\\
    \gll \textit{Eða} \textit{hvat} \textit{sé} \textit{ek} \textit{þar} \textit{upp} \textit{kom-a,} {\textit{\textbf{blik-a}} \iw{ON}{\textit{blika}}} \textit{þar} \textit{eigi} \textit{skild-ir} \textit{við?}\\
      but what see.\textsc{1sg} I there up come-\textsc{inf} \textbf{blika}-\textsc{3sg.prs} there not shield-\textsc{nom.pl} with\\
    \glt ‘But what do I see coming up there, aren’t there shields \textbf{\textit{blika}}?’
\end{exe} 

In the second example, \iwi{ON}{\textit{upp koma}} ‘to come up, appear’ in the first of the two asyndetically connected main clauses conveys an \isi{inchoative} sense of the situation. In my view, this might even speak against a \isi{semelfactive} reading of the second main clause and therefore of \iwi{ON}{\textit{blika}}. As I can see coming up a thunderstorm, but not the lightning that comes with it, I can see coming up a troop of soldiers, but not the flashing of their shields.

Whatever one might think about this question of detail, it seems clear that an unambiguous choice between reading \iwi{ON}{\textit{blika}} as \isi{intensive}-\isi{iterative} ‘to glisten, glint’ or as \isi{inchoative} resp.\ \isi{semelfactive} ‘to flash’ is hardly possible on grounds of the texts that have come down to us. We find ourselves amidst, on the one hand, interpretations that are read into the texts retrospectively and, on the other hand, semantic overlaps that are an integral part of the actual language studied and every language in general. 

\subsection{\textit{brúnn} ‘brown, dark brown; brownish-violet’}
\label{sec:OWNbrown}

As there is no class III weak verb of \iwi{ON}{\textit{bleika}} attested, I shall proceed to the next color adjective on the list, i.e., \iwi{ON}{\textit{brúnn}} ‘brown, dark brown; brownish-violet’.\footnote{For a detailed study on the use of \iwi{ON}{\textit{brúnn}} in Old Icelandic, cf. \citet{Wolf2017}.} With OE, OS, OHG, and OFris.\ \iwi{OFris}{\textit{brūn}}, it can be traced back to PG \iwi{PGmc}{*\textit{brūna}-} (\citetalias[ II 374]{EWA1988}; \citetalias[ 143]{EWGP}). \iw{OE}{\textit{brūn}} \iw{OS}{\textit{brūn}} \iw{OHG}{\textit{brūn}}

A class I verb is not attested in OWN, but in OHG there is \iwi{OHG}{\textit{brūnen}} ‘to make brown-violet, purple; to give colorfulness’ (only in Notker), to be read as /brynen/ with not yet graphically realized \textit{i}-Umlaut, cf. MHG \iwi{MHG}{\textit{briunen}}, German \iwi{German}{\textit{bräunen}}. 

As for class II in OWN, the expected OWN \textit{brúna}*\iw{ON}{\textit{brúna}* (OWN)} ‘to become brown’ is not actually attested. Interestingly, the adjective \iwi{ON}{\textit{brúnaðr}} ‘brown colored’ is attested at least 17 times in the OWN prose texts (according to \citetalias{ONP}). The existence of this adjective presupposes that at some earlier point in time \textit{brúna}*\iw{ON}{\textit{brúna}* (OWN)} must have been a verb that either vanished or just happens not to occur in the texts we know.\footnote{\citet[132, Fn. 30]{Crawford2014} states: “A verb \textit{bruna} \iw{ModNorw}{\textit{brune/-a}} does exist in ModNorw, but it does not have \isi{inchoative} meaning”. This is indeed the case, but the way Crawford puts it is rather misleading. There is a verb \textit{brune/-a},\iw{ModNorw}{\textit{brune/-a}} but according to \href{ordbønene.no}{ordbønene.no}, it has \isi{factitive} semantics ‘gjøre brun’.}

\subsection{\textit{grár} ‘grey; hostile’}
\label{sec:OWNgrey}

Derivations of OWN \iwi{ON}{\textit{grár}} ‘grey’, also ‘hostile’,\footnote{For a detailed study on the use of \iwi{ON}{\textit{grár}} in Old Icelandic, cf. \citet{Wolf2009}.} are comparably sparse. Because of OHG \iwi{OHG}{\textit{grāo}}, OS \iwi{OS}{-\textit{grē}} (only in \iwi{OS}{\textit{appulgrē}} ‘grey spotted, piebald [about a horse]’), Langob. \iwi{Langob}{*\textit{grāus}}, Goth. \iwi{Goth}{*\textit{grēwa}}, we know that there must have been PG \iwi{PGmc}{*\textit{grēwa}-} (with *\textit{grēw-ija}- \iw{PGmc}{*\textit{-ija}}> OE \iwi{OE}{\textit{grǣg}}), another \textit{wa}-adjective just like the words for ‘blue’ and ‘yellow’ (\citetalias[ IV 591]{EWA1988}; \citetalias[ 259]{EWGP}; \citealt[265--266]{Luhr2000}; \citetalias[ 189]{Kroonen2013}).

No class I verb is attested to this adjective. OHG \iwi{OHG}{\textit{grāwēn}} is a class III verb that surfaces only in the glosses. Typical for OHG verbs of this class, it vacillates between \isi{inchoative} and \isi{durative} semantics and therefore would have to be translated as ‘to become grey’ and/or ‘to be grey; shimmer’. 

In OWN, again, there is a class II verb attested only here, namely \textit{grána}\iw{ON}{\textit{grána} (OWN)}. According to the dictionaries, it shows predominantly metaphoric semantics:

\begin{itemize}
\item \citetalias{ONP}: ‘to develop badly (between someone), take a serious turn’
\item \citetalias{Cleasby1874}: ‘to become grey, metaph. to be coarse and spiteful’
\item \citetalias{WAP}: ‘(eigentlich grau werden) unfreundlich werden’
\item \citetalias{OgnS}: ‘blive graa, antage en graa Farve' (‘to become grey, take on a grey color’)
\item \citetalias{LP}: ‘fade away’
\end{itemize}

The semantics in the \citetalias{ONP} are given based on the following example:
 
\begin{exe} 
    \ex Sturlunga saga 20,7 \label{ex:grana}\\
    \gll \textit{Ok} \textit{af} \textit{þess-um} \textit{ákǫst-um} \textit{tek-r} \textit{heldr} \textit{at} {\textit{\textbf{gran-a}} \iw{ON}{\textit{grána} (OWN)}} \textit{gaman-it,} \textit{ok} \textit{kom-a} \textit{kviðling-ar} \textit{við} \textit{svá}\\
      and from this-\textsc{dat.pl} provocation-\textsc{dat.pl} begin-\textsc{3sg.prs} rather to \textbf{take.a.serious.turn}\textsc{-inf} pleasure.\textsc{nom.sg-det} and come-\textsc{inf} {mocking.verse-\textsc{nom.pl}} with so\\
    \glt ‘... and because of these provocations the pleasure began to rather \textbf{take a serious turn} and mocking verses were added.’
\end{exe} 

\subsection{\textit{grǿnn} ‘green’}
\label{sec:OWNgreen}

Somewhat surprisingly, only one verbal derivation\is{derivation!verbal} to OWN \iwi{ON}{\textit{grǿnn}} ‘green’ is reflected in the texts. The adjective itself has parallels in OE, OFris. \iwi{OFris}{\textit{grēne}},\iw{OE}{\textit{grēne}} OS \iwi{OS}{\textit{grōni}}, OHG \iwi{OHG}{\textit{gruoni}} and can therefore be said to continue an older PG \iwi{PGmc}{*\textit{grōni}-} (EWA IV: 667; \citetalias[ 260--261]{EWGP}; \citetalias[ 191--192]{Kroonen2013}).

It is the second time that we encounter a class I verb: \textit{grøna}. According to \citetalias[ s.v.\ \textit{grǿna}]{ONP}, it is attested only once, as medio-passive \textit{grønask}.\iw{ON}{\textit{grønask} (\textit{grøna})} It surfaces in the Flórents saga as a 3sg.\ preterit form <græniz> with the ending -\textit{iz} from earlier (but in this particular case unattested) \iwi{ON}{-\textit{isk}} (\citetalias[ §544,2]{AnGr}). Perhaps because it is attested only once in the prose texts (and not at all in the poetic texts), \citetalias{ONP} and \citetalias{Cleasby1874} agree in their translation, while \citetalias{OgnS} does not list the verb:

\begin{itemize}
\item \citetalias{ONP} and \citetalias{Cleasby1874}: ‘blive grøn’ (‘to become green’)
\end{itemize}

Mindful of the fact that we are dealing here with a medio-passive class I verb, we could expect a reflexive \isi{factitive} meaning ‘to make oneself green’. But who or what is capable of doing such a thing? The solution becomes apparent if we have a look at the actual paragraph: 
 
\begin{exe}
    \ex Flórents saga 197,53 \label{ex:grona}\\
    \gll {...} \textit{þa} \textit{leinga-z} \textit{dag-ar,} \textit{ok} \textit{batna-r} \textit{vedratta} {\textit{\textbf{græni-z}}} {\textit{jord} ...}\\
      {} then lengthen-\textsc{3sg.prs.refl} day-\textsc{nom.pl} and improve-\textsc{3sg.prs} weather.\textsc{nom.sg} \textbf{make.green}-\textsc{3sg.prs.refl} earth.\textsc{nom.sg} {}\\
    \glt ‘... the days get longer, and the weather improves, the earth \textbf{makes itself green} (= becomes green) ...’ \iw{ON}{\textit{grønask} (\textit{grøna})}
\end{exe}

If one believes that the earth is referred to as an animate entity here and is therefore not only the subject of the clause, but also the agent, then the use of \textit{græniz} \iw{ON}{\textit{grønask} (\textit{grøna})}   would be strictly reflexive. If, on the other hand, this is not the case, it would point to \isi{anticausative} semantics, since \iwi{ON}{\textit{jord}} ‘earth’ is still the subject of the clause, but not the agent. As the patient, becoming green is something that happens to it.\footnote{This matter cannot be decided on the basis of the text. There are references in Old Norse sources (especially the translation of the Elucidarius, and the Edda) to the equation common in the Middle Ages of a man as a microcosm with the macrocosm, whereby the four major elements earth, water, air and fire correspond to human flesh, blood, breath and body heat, cf. \citet[94--95]{Simek1990} and \textit{passim}.}

Class III verbs are attested only for OE \iwi{OE}{\textit{grēnian}} ‘to become green (about foliage); flourish’, and OHG \iwi{OHG}{\textit{gruonēn}} ‘to become / to be green, blossom, flourish’ again with the vacillation between \isi{inchoative}/\isi{ingressive} and \isi{durative} meaning so characteristic for OHG class III verbs. 

\subsection{\textit{hvítr} ‘white’}
\label{sec:OWNwhite}

\largerpage

Continuations of PG \iwi{PGmc}{*\textit{hʷeita}-} ‘white’ are widely attested. Next to OWN \iwi{ON}{\textit{hvítr}} we have West Germanic cognates OE \iwi{OE}{\textit{hwīt}}, OHG \iwi{OHG}{(\textit{h})\textit{wīz}}, OFris.\ \iwi{OFris}{(\textit{h})\textit{wīt}}, \textit{wit}, and even Goth.\ \iwi{Goth}{\textit{hveits}} (\citetalias[ 316--317]{EWGP}; \citetalias[ 267]{Kroonen2013}).

A Proto-Germanic \isi{factitive} \iwi{PGmc}{*\textit{hʷeitijan}-} can be assumed based on OE \iwi{OE}{\textit{hwītan}} ‘to make white, to whiten; to polish’, OHG \iwi{OHG}{(\textit{h})\textit{wizen}} ‘\textit{idem}; also: to wash white’, and Goth. \iwi{Goth}{\textit{gahweitjan}} ‘to make white, to whiten’. It is, however, not attested in OWN.\footnote{An online search in the \citetalias{ONP} for the definitions ‘to make white’ or ‘whiten’ does not yield a single result.}

According to the findings so far, we have yet another class II \textit{nan}-verb \textit{hvítna}\iw{ON}{\textit{hvítna} (OWN)} that is restricted to OWN. Again, the dictionaries agree on the meaning:

\begin{itemize}
\item \citetalias{ONP}: ‘blive hvid’ (‘to become white’)
\item \citetalias{Cleasby1874}: ‘to become white’
\item \citetalias{WAP}: ‘weiß werden’
\item \citetalias{OgnS}: ‘blive hvid’ (‘to become white’)
\end{itemize}
 
\begin{exe} 
    \ex Snorra Edda 49,22 \label{ex:hvitna}\\
    \gll \textit{hann} \textit{herþ-i} \textit{hendrn-ar} \textit{at} \textit{hamar-skapt-i-nv,} \textit{sv} \textit{at} {\textit{\textbf{hvitnv-þv}} \iw{ON}{\textit{hvítna} (OWN)}} \textit{knv-ar-nir}\\
     he {clasp.firmly-\textsc{3sg.pret}} hand-\textsc{nom.pl} at hammer-shaft-\textsc{det-dat.sg} so that \textbf{become.white}-\textsc{3pl.pret} knuckle-\textsc{nom.pl-det}\\
    \glt ‘He clasped the hammer shaft so firmly with his hands that his knuckles \textbf{became white}.’
\end{exe} 

Interestingly, the dictionaries contradict one another with respect to verbs attested in OE \iwi{OE}{\textit{hwītian}} (class II) and OHG \iwi{OHG}{\textit{wīzēn}} (class III). According to \citetalias[ 273]{Vries1977}, neither of them is \isi{inchoative}. According to \citetalias[ s.v. \textit{hwítian}]{BT}, however, the English verb vacillates between ‘to be / to become white’. The same is true for the OHG verb. While \citet[397]{Schutzeichel2012} translates it as \isi{inchoative} ‘weiß werden’, Heidermanns (\citetalias[ 317]{EWGP}) suggests ‘weiß werden, sein’. This, again, illustrates that, when it comes to semantics, working with dictionaries can be a venture that should be approached with caution. 

\subsection{\textit{rauðr} and \textit{rjóðr} ‘red’}
\label{sec:OWNred}

We shall now deal with derivations of the words for ‘red’. In addition to OWN \iwi{ON}{\textit{rauðr}}, a second lexeme \iwi{ON}{\textit{rjóðr}} is attested in the texts as well. While \iwi{ON}{\textit{rauðr}} is the common word (at least 248 instances according to \citetalias{ONP}, and 112 according to \citetalias{LP}), \iwi{ON}{\textit{rjóðr}} is “extremely uncommon [and] used only for the color of human faces” \citep[120]{Crawford2014}. 

OWN \iwi{ON}{\textit{rauðr}} can formally be traced back to PG \iwi{PGmc}{*\textit{rauda}-} ‘red’, as shown by OE \iwi{OE}{\textit{rēad}}, OFris. \iwi{OFris}{\textit{rād}}, OS \iwi{OS}{\textit{rōd}}, OHG \iwi{OHG}{\textit{rōt}}, and Goth. \iwi{Goth}{\textit{rauþs}}. OWN \iwi{ON}{\textit{rjóðr}} on the other hand belongs to the ablauting PG \iwi{PGmc}{*\textit{reuda}-} ‘red’, in turn shown by OE \iwi{OE}{\textit{rēod}}, and Goth. \iwi{Goth}{\textit{gariuþs}} ‘honorable’ (< ‘shameful’ < ‘with a red face’), \citetalias[ VII 653--654]{EWA1988}; \citetalias[ 438--439, 448]{EWGP}; \citet[71]{Luhr2000}.\footnote{It has been suggested (\citetalias[ 448]{EWGP}), that PG \iwi{PGmc}{*\textit{reuda}-} ‘red’ is a Germanic innovation to the present stem of a strong verb PG \iwi{PGmc}{*\textit{reuda}-} ‘to make red’, but Ved. \iwi{Sanskrit}{\textit{róhita}-} ‘red (adj.); red horse (subs.)’, Av. \iwi{Avestan}{\textit{raoδita}-} ‘red’, and Lat. \iwi{Latin}{\textit{rūbidus}} < \iwi{PIE}{*\textit{h\textsubscript{1}re/ou̯dʰ-i}-} (EWA VII: 655) show that there was an old \textit{e}/\textit{o}-ablauting adjective of which both forms survived in Germanic and were therefore directly ancestral to \iwi{PGmc}{*\textit{reuda}-} as well as \iwi{PGmc}{*\textit{rauda}-}. Moreover, there must have been a zero-grade \textit{ró}-derivative \iwi{PIE}{*\textit{h\textsubscript{1}rudʰ-ró}-} (> Latin \iwi{Latin}{\textit{ruber}}, Greek \iwi{AG}{ἐρυθρός} ‘red’) that gave rise to PG \iwi{PGmc}{*\textit{rudra}-}. As this only survives as an \textit{ōn}-stem feminine noun OWN \iwi{ON}{\textit{roðra}} ‘blood’, it can be disregarded here.}

If we now take a closer look at the weak verbs, we find \textit{roðna}\iw{ON}{\textit{roðna} (OWN)} and \iwi{ON}{\textit{roða}} (both class II). The first one is translated by the dictionaries as having \isi{anticausative}-\isi{inchoative} semantics:

\begin{itemize}
\item \citetalias{ONP}: –
\item \citetalias{Cleasby1874}: ‘to redden, become red (of the face), to blush’
\item \citetalias{WAP}: ‘rot werden, erröten’
\item \citetalias{OgnS}: ‘rødme, blive rød; blive rød i Kinderne’ (‘to blush, become red; to turn red in the cheeks’)
\end{itemize}

An example for the use of \textit{roðna}\iw{ON}{\textit{roðna} (OWN)} (as \textsc{3sg.pret} \textit{rodnaþi} ‘(she) blushed’) can be seen below in (\ref{ex:roda2}). 

Less clear, on the other hand, is the derivational history of this word. No cognates in the other Germanic languages suggest that it might be an OWN innovation. According to \citet[100]{Wessen1965}, all it takes for a word \textit{roðna}\iw{ON}{\textit{roðna} (OWN)} to be built is the need for finite verb forms with the meaning ‘to become red’ and a sample group (“Mönstergrupp”) that can serve as a model. As we have seen, the other OWN weak \textit{nan}-verbs formed to color adjectives have no correspondences in the other Germanic languages. It can thus be assumed that the pattern was productive\is{productivity} in OWN.

\citet[379]{Seebold1970}, however, states that \textit{roðna}\iw{ON}{\textit{roðna} (OWN)} ‘rot werden, sich röten’ is derived from the strong verb PG \iwi{PGmc}{*\textit{reuda}-} ‘röten’. Heidermanns’ etymological dictionary of the Germanic primary adjectives\is{derivation!deradical} does not contain \textit{nō}-derivatives neither s.v. \textit{rauda}- ‘rot’ (\citetalias[ 438--439]{EWGP}) nor s.v. \textit{reuda}- ‘gerötet’ (\citetalias[ 448]{EWGP}). He therefore seems to assume that OWN \textit{roðna}\iw{ON}{\textit{roðna} (OWN)} is of deverbal origin as well.\is{deverbal/deverbative!verb}

If we were to assume that OWN \textit{roðna}\iw{ON}{\textit{roðna} (OWN)} continues a \textit{voreinzelsprachliche} formation, we could not decide if it was deverbal\is{deverbal/deverbative!verb} from \iwi{PGmc}{*\textit{reuda}-} ‘redden’ or deadjectival \is{deadjectival!verb} from either \iwi{PGmc}{*\textit{rauda}-} ‘red’ or \iwi{PGmc}{*\textit{reuda}-} ‘reddened’ on formal or semantic grounds. On the formal side, all three possibilities would require a zero-grade root and the subsequent development *\textit{rud-nō}- > *\textit{rud-na}- > \textit{roðna}\iw{ON}{\textit{roðna} (OWN)}. 

On the semantic side, if we recall the \isi{anticausative}-\isi{inchoative} function of the \textit{nō}-formations, a deverbal\is{deverbal/deverbative!verb} formation to \iwi{ON}{*\textit{reuda}-} would yield ‘to become reddened’,\footnote{Cf.\ Goth.\ \textit{us-brukan} \iw{Goth}{\textit{usbrukan}} ‘to become broken off’ : \iwi{Goth}{\textit{brikan}} ‘to break’, OWN. \iwi{ON}{\textit{sofna}} ‘to fall asleep’ : \iwi{ON}{\textit{sofa}} ‘to sleep’, see \citet[253]{KraheMeid1969} for further examples.} and a deadjectival \is{deadjectival!verb} formation would yield ‘to become red’ or ‘to become reddened’.

Given that there are no cognates in the other Germanic languages and that former \textit{nō}-formations were productive\is{productivity} in OWN, the assumption that \textit{roðna}\iw{ON}{\textit{roðna} (OWN)} is indeed an innovation seems to be the most plausible.

In addition to \textit{roðna}\iw{ON}{\textit{roðna} (OWN)} there is, as stated above, another class II weak verb: \iwi{ON}{\textit{roða}}. No translation is given in the \citetalias{ONP} yet, but the other dictionaries do not contradict one another regarding the meaning of the verb:

\begin{itemize}
\item \citetalias{Cleasby1874}: ‘to gleam red’
\item \citetalias{WAP}: ‘roten Schein geben, rot leuchten’
\item \citetalias{OgnS}: ‘kaste et rødt Skin, en rød Farve fra sig eller over noget’ (‘to cast a red light, a red color from oneself or at something’)
\end{itemize}

In the Germanic etymological dictionaries, \iwi{ON}{\textit{roða}} is usually reconstructed as a class III verb \iwi{PGmc}{*\textit{rudēn}-} (> OHG \iwi{OHG}{\textit{rotēn}} ‘to be or become red’, \citetalias[ 450]{Vries1977}; \citetalias[ 416]{Kroonen2013}; \citetalias[ VII 658]{EWA1988}). While this is formally possible, the attested Old West Norse forms, however, synchronically show class II inflection. According to \citetalias{ONP}, the verb is attested at least twice in the Old West Norse prose literature. There are, however, no attestations in the poetic texts. Two of the attestations are given in examples (\ref{ex:roda1})-(\ref{ex:roda2}) below.

 
\begin{exe}
    \ex Sverris saga 92,29 \label{ex:roda1}\\
    \gll \textit{Sva} \textit{var} \textit{at} \textit{sia} \textit{i} \textit{fiall-it} \textit{upp} \textit{sem} \textit{i} \textit{log-a} \textit{sęi,} \textit{er} {\textit{\textbf{roþa-ði}} \iw{ON}{\textit{roða}}} \textit{sciolld-unum}\\
      so was to look.\textsc{inf} at rock.\textsc{nom.sg-det} up like  in fire-\textsc{acc.sg} look.\textsc{3sg.pst.sbjv} that \textbf{gleam.red}-\textsc{3sg.pret} shield-\textsc{dat.pl}\\
    \glt   ‘Looking up at the rock that \textbf{gleamed red} from the shields was like looking into a fire.’

    \ex Karlamagnúss saga 199,16 \label{ex:roda2}\\
    \gll {...} \textit{Kong-s} \textit{dottir} {\textit{rodna-þi} \iw{ON}{\textit{roðna} (OWN)}} \textit{miǫg} \textit{vid} \textit{þess-a} \textit{sǫgu} \textit{ok} \textit{vard} \textit{þvi} \textit{likuz} \textit{sem} \textit{þa} \textit{er} {\textit{\textbf{roda-r}} \iw{ON}{\textit{roða}}} \textit{fyrir} \textit{vpprennandi} \textit{solu} \textit{i} \textit{hinv} \textit{fegursta} \textit{heidi}\\
    {} king-\textsc{gen.sg} daughter.\textsc{nom.sg} blush-\textsc{3sg.pret} {very.much} at this-\textsc{acc.sg} story.\textsc{acc.sg} and become.\textsc{3sg.pret} therefore like as then when \textbf{gleam.red}-\textsc{3sg.prs} while rise.\textsc{ptcp.prs.dat.sg} sun.\textsc{dat.sg} in the.\textsc{art.dat.sg} beautiful.\textsc{superl} clear.sky.\textsc{dat.sg}\\
    \glt ‘The king's daughter blushed very much at this story and (she) therefore then became as when it \textbf{gleams red} while the rising of the sun in its most beautiful clear sky.’
\end{exe} 

The two forms in question are \textit{roþaði} \iw{ON}{\textit{roða}} ‘(the rock) gleamed red’ (\textsc{3sg.pret}) in (\ref{ex:roda1}) and the impersonal \textit{rodar} \iw{ON}{\textit{roða}} ‘it gleamed red’ (\textsc{3sg.prs}) in (\ref{ex:roda2}). In Old West Norse, only the class II weak verbs show a vowel -\textit{a}- in the middle syllable of the present and preterit forms, whereas the class III verbs have no such vowel in the preterit nor -\textit{i}- in the present, cf.\ \textit{vakti} (\textsc{3sg.pret}), \textit{vakir} (\textsc{3sg.prs}) from \iwi{ON}{\textit{vaka}} ‘to be awake’.

Formally, \iwi{ON}{\textit{roða}} could continue both a class II and class III verb, but the \isi{durative}/\isi{essive} semantics rather point to a class III verb that secondarily has adopted class II inflection.

\subsection{\textit{svartr} ‘black, dark’}
\label{sec:OWNblack}

Let us finally have a look at the verbal derivations\is{derivation!verbal} of the word for ‘black, dark’. OWN \iwi{ON}{\textit{svartr}}, Goth. \iwi{Goth}{\textit{swarts}}, OE \iwi{OE}{\textit{sweart}}, OHG \iwi{OHG}{\textit{swarz}} indicate that it continues PG \iwi{PGmc}{*\textit{swarta}-} (\citetalias[ 574]{EWGP}).

While an already PG \isi{factitive} \iwi{PGmc}{*\textit{swartijan}-} ‘to make black/dark, blacken/darken’ can be reconstructed on the basis of Goth. \iwi{Goth}{\textit{swartjan}*}, OHG \iwi{OHG}{\textit{swerzen}} as well as ODan. \iwi{ODan}{\textit{swertæ}} and OSwed. \iwi{OSwed}{\textit{sværta}}, the West Norse word is attested no earlier than in ModNorw \iwi{ModNorw}{\textit{sverte}} (\citealt[1222]{BjorvandLindeman2019}).

In contrast, class II verbs are attested in OWN only. Here we have another \textit{nan}-verb with zero-grade root, i.e., \textit{sortna}\iw{ON}{\textit{sortna} (OWN)}. Again, the \citetalias{ONP} does not offer a translation, but the other dictionaries agree on \isi{anticausative}-\isi{inchoative} semantics:

\begin{itemize}
\item \citetalias{Cleasby1874}: ‘to grow black’
\item \citetalias{WAP}: ‘schwarz, dunkel werden’
\item OngS: ‘sortne, antage en sort Farve’ (‘to become black / blacken, take on a black color’)
\end{itemize}
 
\begin{exe} 
    \ex Cecilíu saga 284,4 \label{ex:sortna}\\
    \gll {...} \textit{dagr-inn} \textit{fly-di,} \textit{ok} \textit{sol} {\textbf{\textit{sortna-di,}} \iw{ON}{\textit{sortna} (OWN)}} \textit{ok} \textit{ger-di} \textit{myrk-t} \textit{i} \textit{heim-inum} \textit{aull-um}\\
     {} day.\textsc{nom.sg-det} leave-\textsc{3sg.pret} and sun.\textsc{nom.sg} \textbf{blacken}-\textsc{3sg.pret} and make-\textsc{3sg.pret} dark-\textsc{n} in home-\textsc{dat.pl.det} all-\textsc{dat.pl}\\
    \glt ‘... the day left and the sun \textbf{blackened} and made all homes dark.’
\end{exe} 

More interesting, however, is \iwi{ON}{\textit{sorta}}. According to \citetalias{ONP} and \citetalias{LP} this class II verb is sparsely attested in Old Norse prose, but at least four times there and once in the poetic texts. While \citetalias{WAP} and \citetalias{ONP}, again, offer no translation, a look in the other dictionaries does not reveal contradictions:

\begin{itemize}
\item \citetalias{Cleasby1874}: ‘to dye black’
\item OsgN: ‘formørke’ (‘to turn sth. black/dark; to eclipse’) (tr.)
\item \citetalias{LP}: ‘to tar’
\end{itemize}

I will still discuss this word here because class II verbs are not usually associated with what seems like \isi{causative}/\isi{factitive} semantics.

From the four instances that are listed in the \citetalias{ONP}, I will discuss two instructive examples here. The first one in example (\ref{ex:sorta1}) describes people living in Ethiopia:
 
\begin{exe}
    \ex Stjórn 96,18 \label{ex:sorta1}\\
    \gll {...} \textit{hueria} \textit{er} \textit{sol-in} \textit{brenn-ir} \textit{ok} {\textit{\textbf{sorta-r}} \iw{ON}{\textit{sorta}}}\\
     {}  which.\textsc{acc.pl} who sun.\textsc{nom.sg-det} burn-\textsc{3sg.prs} and \textbf{sorta}-\textsc{3sg.prs}\\
    \glt ‘...which the sun burns and \textbf{\textit{sortar}}.’
\end{exe} 

The first component of the binomial \iwi{ON}{\textit{brennir}} ‘burns’ (\textsc{3sg.prs}) (< PG \iwi{PGmc}{*\textit{brannjan}-} ‘to make burn’ tr.), which in OWN falls together with its derivational basis PG \iwi{PGmc}{*\textit{brinnan}-} ‘burn’ (intr.; \citetalias[ 56]{Vries1977}),\footnote{Something similar has happened or is happening in most Germanic languages. For example, until Early New High German the strong verb \iwi{ENHG}{\textit{brinnen}} was distinct from its derivative \iwi{ENHG}{\textit{brennen}}. The former is in use today only in some Upper German dialects and is being replaced by the latter (\citetalias[ II 342]{EWA1988}).} conveys \isi{causative} semantics. This seems to be true also for the second one, \textit{sortar} \iw{ON}{\textit{sorta}} ‘makes black, blackens’ (\textsc{3sg.prs}) – at least at first glance. However, this would be rather unexpected since it is the class I verbs that are known to be \isi{causative}s/\isi{factitive}s, an observation that is confirmed by the material discussed here as well.

We might assume based on the aforementioned forms ODan. \iwi{ODan}{\textit{swertæ}}, OSwed. \iwi{OSwed}{\textit{sværta}}, and ModNorw.\ \iwi{ModNorw}{\textit{sverte}} that there was in fact once an OWN \isi{factitive} \iwi{ON}{*\textit{sverta}} ‘to make black/dark, blacken’ that is coincidentally not attested in the texts. If \iwi{ON}{\textit{sorta}} also had \isi{factitive} semantics, as suggested by \citetalias{Cleasby1874}, OsgN and \citetalias{LP}, we would have to assume that these two forms were in \isi{competition} with one another and that \textit{sorta} did not emerge victorious. This is far from impossible. However, this issue cannot be solved with the material the dictionaries provide, but only by means of a thorough investigation of the actual texts in their local and temporal context.

At second glance, however, another option seems more promising. We can see that Heidermanns connects \iwi{ON}{\textit{sorta}} to the zero-grad nouns OWN \iwi{ON}{\textit{sorti}} ‘blackness, darkness’ and OWN \iwi{ON}{\textit{sorta}} ‘black color’ (\citetalias[ 574]{EWGP}). Since Germanic denominal \is{denominal!verb} class II verbs usually exhibit the same ablaut grade in the root as their derivational bases (\citealt{KraheMeid1969}: 238–240), and since no strong verb that could have functioned as derivational base exists, assuming a denominal \is{denominal!verb} formation for \iwi{ON}{\textit{sorta}} is formally appealing. Reconsidering the formal aspects of denominal \textit{ōn}-formations \is{denominal!verb} enables us to better understand the role \textit{sortar} \iw{ON}{\textit{sorta}} plays in Stjórn 96,18 in example (\ref{ex:sorta1}). A case in point is OHG \iwi{OHG}{\textit{salbōn}}, OS \iwi{OS}{\textit{salvon}}, OE \iwi{OE}{\textit{sealfian}} < West Germanic \iwi{PGmc}{*\textit{salbōn}-}, a derivative of \iwi{PGmc}{*\textit{salbō}-} ‘ointment’. The verb is best paraphrased as ‘to provide with ointment’. The same is true for Goth. \iwi{Goth}{\textit{gapaidōn}}, a derivative of \iwi{Goth}{\textit{paida}} ‘garment’, best paraphrased as ‘provide with garment’ (\citealt[239]{KraheMeid1969}) or OWN \iwi{ON}{\textit{dúka}}, a derivative of \iwi{ON}{\textit{dúkr}} ‘cloth, veil, scarf’ and therefore ‘to provide with a \iwi{ON}{\textit{dúkr}}’ \citep[178]{Iversen1923}. These class II verbs have been termed ‘ornative’ (recently \citealt[250]{Filatkina2018}; \citealt[376--378]{Schaefer1984}), and this is key to understanding \iwi{ON}{\textit{sorta}} as well. For example, in contrast to OWN \iwi{ON}{\textit{dǿma}} ‘to judge, pronounce a judgement’, a class I derivative of \iwi{ON}{\textit{dómr}} ‘judgement, verdict’, ornative verbs do not denote the realization or materialization of what is denoted in the derivational basis, but that something or someone is provided with the feature or property that the derivational basis has. This means that \iwi{ON}{\textit{sorta}} conveys a meaning ‘to provide with blackness’, rather than a meaning ‘to make black, to blacken’. As the difference is not as clear cut as, for example, in \iwi{PGmc}{*\textit{salbōn}-} (‘to provide with ointment’ vs. hypothetical *‘to make/turn into ointment’), constellations like these might have contributed to the collapse of the morphosyntactic system of the Germanic weak verbs. 

Let us now see if this interpretation can hold up in the light of another example:
 
\begin{exe} 
    \ex Konungs skuggsjá 60,29 \label{ex:sorta2}\\
    \gll {...} \textit{En} \textit{þo} \textit{ero} \textit{pannzar-ar} \textit{hofuð} \textit{vapn} \textit{til} \textit{lifð-ar} \textit{a} \textit{skip-um} \textit{gorv-ir} \textit{af} \textit{blaut-um} \textit{lerept-um} \textit{oc} \textit{vel} {\textit{\textbf{sortu-dum}}\iw{ON}{\textit{sorta}}}, \textit{goð-ar} \textit{hialm-ar} \textit{oc} \textit{hang-anða} \textit{stalhufur.}\\
   {} and but be.\textsc{3pl.pret} shield\textsc{-nom.pl} main weapon.\textsc{nom.sg} for protection\textsc{-gen.sg} on ship-\textsc{dat.pl} make-\textsc{ptcp.prf.nom.pl} of soft-\textsc{dat.pl} canvas-\textsc{dat.pl} and well \textbf{sorta}-\textsc{ptcp.prf.dat.pl} good-\textsc{nom.pl} helmet-\textsc{nom.pl} and hang-\textsc{ptcp.prs.nom.pl} {steel.hood-\textsc{nom.pl}}\\
    \glt ‘But shields are the main protective weapons on the ships, made of soft and well-\textbf{blackened} canvas, good helmets and low-reaching steel hoods.’
\end{exe} 

It is a reasonable assumption that the canvas is the one that is ‘provided with blackness’ and that the \isi{participle} consequently means ‘blackened’.

\section{Conclusion}
\label{sec:OWNconclusion}

Let us now wrap up the most important findings. Class I verbal derivatives\is{derivation!verbal} are attested for three out of the eight BCT \is{Basic Color Term (BCT)} that I have discussed here: OWN \iwi{ON}{\textit{bleikja}}, \iw{ON}{\textit{grønask} (\textit{grøna})} \textit{grønask}, and – in lack of an OWN cognate – ODan. \iwi{ODan}{\textit{swertæ}} resp. OSwed. \iwi{OSwed}{\textit{sværta}}. All of them are at least of Proto-Germanic origin and, according to the dictionaries, express \isi{factitive} semantics.

The situation is the opposite with regard to the class II \textit{nan}-verbs. For OWN \textit{blána}\iw{ON}{\textit{blána} (OWN)}, \textit{grána}\iw{ON}{\textit{grána} (OWN)}, \textit{hvítna}\iw{ON}{\textit{hvítna} (OWN)}, \textit{roðna}\iw{ON}{\textit{roðna} (OWN)}, \textit{sortna}\iw{ON}{\textit{sortna} (OWN)}, and \textit{brúna}*\iw{ON}{\textit{brúna}* (OWN)} (which can be assumed to have existed based on the adjective \iwi{ON}{\textit{brúnaðr}} ‘brown colored’), no Germanic cognates are attested. This group was highly productive\is{productivity} in OWN, and it is doubtful that any of these words are old.\footnote{According to \citet[114--115]{Gorbachov2007} it is unclear how and if OWN \textit{hvítna}\iw{ON}{\textit{hvítna} (OWN)} is related to Lith.\ \textit{šviñta} (Inf.\ \textit{švìsti})\iw{Lith}{\textit{švìsti, šviñta, švìto}, 1sg.\ prs.\ \textit{švintù}} ‘become light, start shining’. The long vowel in the OWN form could be explained as the result of leveling with its base \iwi{ON}{\textit{hvítr}}, but the consonants do not match. No mention is made of this connection in \citetalias[ s.v.\ \textit{švìsti}]{ALEW}. \citet[116]{Gorbachov2007} also shows that OWN \textit{roðna}\iw{ON}{\textit{roðna} (OWN)} is a perfect formal match to OIr.\ \iwi{MIr}{\textit{roindid}} [sic!] ‘reddens, dyes (with blood)’ and Ved. \iwi{Sanskrit}{\textit{ruṇáddhi}} ‘restrains, keeps back’, but the latter may very well represent a different PIE verbal category altogether since they have a different (\isi{transitive}-terminative) function. As Stefan Höfler (p.c.\ 17.01.25) has furthermore informed me, \textit{roindid} is rather Middle Irish with analogous palatalization of the cluster -\textit{nd}-. For the OIr.\ form see \textsc{Höfler \& Stifter}, this volume, Section 2.2.} All of them express \isi{anticausative}-\isi{inchoative} semantics.\footnote{\citet[346]{Ottosson2013}: ``At least the bulk of the deadjectival \is{deadjectival!verb} [verbs] seems to have denoted types of events that were typically agentless [...].”}

The remaining three verbs \iwi{ON}{\textit{blika}}, \iwi{ON}{\textit{roða}}, and \iwi{ON}{\textit{sorta}}, despite being class II verbs, differ in terms of meaning: The meaning of \iwi{ON}{\textit{blika}} and \iwi{ON}{\textit{sorta}} is disputed according to the dictionaries consulted in this article. Even when looking at actual attestations, an unambiguous determination in terms of semantics is not possible for \iwi{ON}{\textit{blika}}, at least not based on the contexts I investigated. They vacillate between \isi{intensive} and \isi{inchoative} (or \isi{semelfactive}) semantics. 

For \iwi{ON}{\textit{sorta}}, the two dictionaries consulted in this study that offer a translation (\citetalias{Cleasby1874} and \citetalias{OgnS}) are not precise enough in determining its semantic value. With this in mind, the translations suggested in \citetalias{Cleasby1874} and \citetalias{OgnS} ‘to dye black, blacken’ should be refined to ‘to provide with blackness’. While this might be considered an artificial differentiation and while it is probably handy to stick to ‘to dye/make black’ in a translation within a wider context, with respect to the derivational history of this word it is preferable to be as precise as possible in order to be able to understand how the semantic change may have come about.

Finally, synchronically \iwi{ON}{\textit{roða}} is a class II verb as well, but its \isi{essive}/\isi{durative} semantics suggest that historically it is an old \textit{ēn}-verb that was integrated into class II in OWN.

None of the weak verbs formed to adjectives for BCT \is{Basic Color Term (BCT)} in Old West Norse is a member of the third class, which has become highly productive\is{productivity} in OHG in forming \isi{inchoative} verbs (e.g., \iwi{OHG}{\textit{fūlēn}} ‘to rot or to decay’, \iwi{OHG}{\textit{altēn}} ‘to become or to grow old’).

We know from corpus studies on the other Old Germanic languages that the word forming verbal suffixes exhibit a functional variety already in the earliest texts (see fn. \ref{FNOWNlit}). Along these lines, \citet{Schwerdt2008} has shown that the Old High German verbal classes are often ambiguous with respect to their semantic roles. This is the case for about 60\% of \textit{jan}-verbs, 50\% of \textit{ōn}-verbs, and 50\% of \textit{ēn}-verbs. This ambiguity, however, is not found within Old West Norse, at least to this degree.\footnote{\citet{Boldtinpress}.} In other words: The system of the Old West Norse weak verbs was intact (although it is clear that it could never have been completely consistent) to a degree that change of state \is{change of state (CoS)} verbs could not function as \isi{stative} verbs and \textit{vice versa} as they can in Old High German, i.e., they either denote ‘to be x' or ‘to become x'. However, if and to which degree this might be an artifact arising from working with dictionaries – which do not describe the semantic differences adequately – rather than corpora, is something that deserves to be fleshed out in future studies.

\section*{Acknowledgements}

This article has benefited a lot from the helpful comments of the reviewers, the editors, and Charles Comerford. I would like to thank them.


\section*{Abbreviations}

\begin{tabularx}{.5\textwidth}{lQ}
Av. & Avestan\\
Goth. & Gothic\\
Langob. & Langobardic\\
Lat. & Latin\\
Lith. & Lithuanian\\
MHG & Middle High German\\
MLG & Middle Low German\\
ModNorw. & Modern Norwegian\\
ODan. & Old Danish\\
OE & Old English\\
\end{tabularx}
\begin{tabularx}{.48\textwidth}{lQ}
OFris. & Old Frisian\\
OHG & Old High German\\
OIr. & Old Irish\\
OS & Old Saxon\\
OSwed. & Old Swedish\\
OWN & Old West Norse\\
PG & Proto-Germanic\\
PIE & Proto-Indo-European\\
PN & Proto-Norse\\
Ved. & Vedic\\
\end{tabularx}




{\sloppy\printbibliography[heading=subbibliography,notkeyword=this]}
\end{document}
